\origpage{254}{ROBERTSON AND SEYMOUR}
\noindent circumstances is the induced tangle respectful? We are also concerned with the parallel question when we add a crosscap.

Let $\Sigma$ be a connected surface, and let $C_1, \dots, C_t$ be mutually disjoint $O$-arcs in $\Sigma$. Let $\Sigma_1, \Sigma_2$ be 2-manifolds with boundary, such that $\Sigma_1 \cup \Sigma_2 = \Sigma$ and
\begin{equation*}
\Sigma_1 \cap \Sigma_2 = bd(\Sigma_1) = bd(\Sigma_2) = C_1 \cup \dots \cup C_t.
\end{equation*}
Let $\hat{\Sigma}_1$ be a surface with $\Sigma_1 \subseteq \hat{\Sigma}_1$, such that each $C_i$ bounds an open disc $D_i \subseteq \hat{\Sigma}_1$, and $D_1, \dots, D_t$ are mutually disjoint, and $\hat{\Sigma}_1 - \Sigma_1 = D_1 \cup \dots \cup D_t$. If $t = 2$, $\hat{\Sigma}_1$ is connected, and $\Sigma_2$ is homeomorphic to $\{(x, y) \colon 1 \le x^2 + y^2 \le 2\}$ we say $\Sigma_2$ is a handle of $\Sigma$. Similarly, if $t = 1$, $\hat{\Sigma}_1$ is connected, and $\Sigma_2$ is homeomorphic to a M\"obius band, we say that $\Sigma_2$ is a crosscap of $\Sigma$. We wish to prove the following two results.

\stmt
(5.1) \textit{Let $\Sigma_2$ be a handle of $\Sigma$, and let $\Sigma_1, \hat{\Sigma}_1$ be defined as above. Let the two components of $\Sigma_1 \cap \Sigma_2$ be $C_1, C_2$. Let $\Gamma$ be a drawing in $\Sigma$ such that $C_1, C_2$ are $\Gamma$-normal, and}
\begin{equation*}
|V(\Gamma) \cap C_1| = |V(\Gamma) \cap C_2| = \theta,
\end{equation*}
\textit{and such that $\Gamma \cap \Sigma_2$ has $\theta$ components, each a path with one end in $C_1$ and the other in $C_2$. Suppose that every $\Gamma$-normal line $F \subseteq \Sigma_1$ with one end in $C_1$ and one end in $C_2$ satisfies $|F \cap V(\Gamma)| \ge \theta$. Let $\mathscr{T}_1$ be a tangle in $\Gamma \cap \hat{\Sigma}_1$, respectful towards $\hat{\Sigma}_1$, of order $\theta$, such that $V(\Gamma) \cap C_1, V(\Gamma) \cap C_2$ are free with respect to $\mathscr{T}_1$. Then the tangle in $\Gamma$ induced by $\mathscr{T}_1$ is respectful towards $\Sigma$.}

\stmt
(5.2) \textit{Let $\Sigma_2$ be a crosscap of $\Sigma$, and let $\Sigma_1, \hat{\Sigma}_1$ be defined as above. Let $\Gamma$ be a drawing in $\Sigma$ such that $\Sigma_1 \cap \Sigma_2$ is $\Gamma$-normal, and $|V(\Gamma) \cap \Sigma_1 \cap \Sigma_2| = 2\theta$, and $\Gamma \cap \Sigma_2$ has $\theta$ components, each a path with ends some opposite pair of vertices in $V(\Gamma) \cap \Sigma_1 \cap \Sigma_2$, in the natural sense. Suppose that $\mathscr{T}_1$ is a tangle in $\Gamma \cap \hat{\Sigma}_1$, respectful towards $\hat{\Sigma}_1$, of order $\theta\fix{{} \ge 2}$,\fixnote{fix:5.2}{As printed, (5.2) fails for $\theta = 1$. Let $\Sigma$ be the projective plane, cut into a closed disc $\Sigma_1$ and a M\"obius band $\Sigma_2$. Let $\Gamma$ have two opposite vertices $u,v$ on $\Sigma_1 \cap \Sigma_2$, a path in $\Sigma_1$ joining them, and a path in $\Sigma_2$ joining them that misses the core circle of the M\"obius band. The tangle of order 1 is respectful towards the sphere $\hat{\Sigma}_1$. But the core circle misses $\Gamma$ and bounds no disc, so the induced tangle is not respectful towards $\Sigma$. For $\theta \ge 2$, (5.2) holds, and its use in Graph Minors XV has $\theta \ge 2$.} and every subset $V(\Gamma) \cap \Sigma_1 \cap \Sigma_2$ of cardinality $\le \theta$ is free with respect to $\mathscr{T}_1$. Then the tangle in $\Gamma$ induced by $\mathscr{T}_1$ is respectful towards $\Sigma$.}

Let us formulate a common generalization of (5.1) and (5.2).

\stmt
(5.3) \textit{Let $\Sigma$ be a connected surface, and let $C_1, \dots, C_t$ be mutually disjoint $O$-arcs in $\Sigma$. Let $\Sigma_1, \Sigma_2$ be 2-manifolds with boundary, such that $\Sigma_1 \cup \Sigma_2 = \Sigma$ and}
\begin{equation*}
\Sigma_1 \cap \Sigma_2 = bd(\Sigma_1) = bd(\Sigma_2) = C_1 \cup \dots \cup C_t.
\end{equation*}
\textit{Let $\hat{\Sigma}_1$ be a surface with $\Sigma_1 \subseteq \hat{\Sigma}_1$, such that each $C_i$ bounds an open disc $D_i \subseteq \hat{\Sigma}_1$, and $D_1, \dots, D_t$ are mutually disjoint, and $\hat{\Sigma}_1 - \Sigma_1 = D_1 \cup \dots \cup D_t$.}
